\Needspace{10\baselineskip}
\centerline{\textsc{New Comet.}}

\medskip
Monatliche Correspondenz. Vol. IX. Pp. 432---435. May 1804.

\medskip
The new comet discovered by Dr. \textsc{Olbers} on 12 March was also discovered in Marseille, at the naval observatory, by \textsc{Pons} on 7 March. On 10 March the astronomer \textsc{Bouvard} of the Paris National Observatory also discovered it. At $15^u\ 57^m$ mean time he found its right ascension to be $220^\circ\ 4'$ and its southern declination $1^\circ\ 41'$. On 11 March at $15^u\ 39^m$ mean time he found a right ascension of $220^\circ\ 10'$ and a northern declination of $2^\circ\ 44'$. On 29 March at $9^u$ mean time he found a right ascension of $218^\circ\ 40'$ and a northern declination of $46^\circ\ 37'$.

\textsc{Messier} also followed this comet in Paris. On 17 March at $11^u\ 34^m\ 18^s$ true time he observed right ascension $220^\circ\ 29'\ 45''$ and northern declination $26^\circ\ 18'\ 24''$.

\Needspace{18\baselineskip}
Dr. \textsc{Olbers} observed this comet on ten days through 1 April. From these observations Dr. \textsc{Gauss} obtained the following parabolic elements of its orbit:

\begin{center}
\begingroup\fontsize{9}{11}\selectfont\setlength{\tabcolsep}{3pt}
\begin{tabular}{lr}
Time of perihelion, 13 February 1804 . . . . . & $14^u\ 49^m\ 51^s$ mean Seeberg time \\
Logarithm of the distance . . . . . . . . & $0.0298575$ \\
\end{tabular}\endgroup
\end{center}
{\fontsize{8}{10}\selectfont
\[\left.\begin{array}{ll}\text{Longitude of perihelion . . . . . . . . . .}&148^\circ\ 44'\ 51''\\\text{Longitude of the ascending node . . . . . .}&176\ 47\ 58\end{array}\right\}\begin{array}{l}\text{both measured from}\\\text{the mean equinox}\end{array}\]}
\begin{center}
\begingroup\fontsize{9}{11}\selectfont\setlength{\tabcolsep}{3pt}
\begin{tabular}{lr}
Inclination of the orbit . . . . . . . . . . . & $56\ 28\ 40$ \\
Motion . . . . . . . . . . . . . . & prograde. \\
\end{tabular}\endgroup
\end{center}

Using these elements, the comparison with all of \textsc{Olbers}’s observations is as follows:

\begin{center}
\begingroup\fontsize{9}{11}\selectfont\setlength{\tabcolsep}{3pt}\renewcommand{\arraystretch}{1.12}
\begin{tabular}{r|r|r|r}
1804 & \shortstack{Mean time\\Bremen} & \multicolumn{2}{c}{\shortstack{Difference between calculation\\and observation}} \\ \cline{3-4}
 & & in R.A. & in decl. \\ \hline
12 March & $12^u\,56^m\,13^s$ & $+9''$ & $0''$ \\
13 & $11\,40\,43$ & $-19$ & $-7$ \\
14 & $12\,22\,26$ & $+22$ & $+92$ \\
15 & $8\,54\,41$ & $-31$ & $-79$ \\
20 & $9\,22\,52$ & $-247 ::$ & $+121 ::$ \\
\end{tabular}\endgroup
\end{center}
\begin{center}
\begingroup\fontsize{9}{11}\selectfont\setlength{\tabcolsep}{3pt}\renewcommand{\arraystretch}{1.12}
\begin{tabular}{r|r|r|r}
1802 & \shortstack{Mean time\\Bremen} & \multicolumn{2}{c}{\shortstack{Difference between calculation\\and observation}} \\ \cline{3-4}
 & & in R.A. & in decl. \\ \hline
22 March & $8^u\,59^m\,13^s$ & $-143''$ & $+85''$ \\
27 & $8\,59\,43$ & $-55$ & $-152$ \\
28 & $8\,28\,2$ & $-27$ & $+32$ \\
29 & $8\,45\,41$ & $-26$ & $+33$ \\
1 April & $9\,1\,52$ & $+28$ & $-1$ \\
\end{tabular}\endgroup
\end{center}

\noindent Here the errors vary so irregularly, especially in declination, that they may reasonably be attributed in large part to the observations. Dr. \textsc{Olbers} himself describes the observation of 20 March as quite doubtful.

In this comparison Dr. \textsc{Gauss} took account of. both nutation and aberration, as well as parallax. The latter is not entirely negligible, particularly for the first observations, when the distance of the com%
\SourceLeaf{253}%
et from Earth was only a quarter of the Sun’s distance. He was therefore all the more willing to take it into account, since it adds almost no work.
For these calculations he used the method explained in the excellent paper printed on p. 385 [p. 94 of this volume] of the present issue. For the constants $A$, $B$, $C$, $a$, $b$, $c$ he found the following values for the direct calculation of right ascension and declination:
\[
A = 268^\circ\ 14'\ 21'', \quad B = 176^\circ\ 29'\ 48'', \quad C = 2^\circ\ 19'\ 58''
\]
\[
a = 87\ \ 19\ \ 38, \quad b = 56\ \ 56\ \ 7, \quad c = 33\ \ 12\ \ 3
\]

\Needspace{22\baselineskip}
\indent The \emph{constants} already refer to the true equator and thus include the nutation, which is to be regarded as constant over so short a time. If $v$ is the true anomaly of the comet, the coordinates $x$, $y$, $z$ are therefore found at once by the following very simple formulas:

\[
\begin{aligned}
x &= \frac{\alpha \sin(v + 240^\circ\ 11'\ 14'')}{\cos\frac{1}{2}v^2}, \\[5pt]
y &= \frac{\ObservedBeta \sin(v + 148^\circ\ 26'\ 41'')}{\cos\frac{1}{2}v^2}, \\[5pt]
z &= \frac{\gamma \sin(v + 334^\circ\ 16'\ 51'')}{\cos\frac{1}{2}v^2}
\end{aligned}
\]

\noindent and where

\[
\log\alpha = 0.029384, \quad \log \ObservedBeta = 9.953130, \quad \log \gamma = 9.768302
\]

\indent The quantities $\alpha$, $\ObservedBeta$, $\gamma$ are here the products of the distance at perihelion and the sine of $a$, $b$, $c$, respectively. The angles $240^\circ\ 11'\ 14''$ and so on are equal to the angles $A$, $B$, $C$ less the distance of the node from perihelion $=\,28^\circ\ 3'\ 7''$. — — —

\begin{flushright}
[\textsc{von Zach}.]
\end{flushright}

\bigskip
